\chapter*{Preface}
\addcontentsline{toc}{chapter}{Preface}
\markboth{Preface}{Preface}

This book is for a student preparing for machine-learning questions in Indian campus placements --- the
on-campus online assessments and technical interviews at IITs, NITs and IIITs --- who wants to understand the
material well enough to derive it, compute it by hand under time pressure, implement it from a blank file, and
survive three ``why?''s in a row. Its companion, \otherbook{Deep Learning}, takes over where this one stops.

\section*{What placements actually ask}
A collection of 49 non-coding ML/DL questions from real company papers (Info Edge, Fujitsu's papers for IISc and
IIT Kharagpur, Qualcomm's for IISc, and EXL) breaks down as: deep learning 22, supervised models 9,
optimisation, loss and regularisation 7, preprocessing, NLP and miscellaneous 6, unsupervised 3, evaluation 1,
probability and statistics 1. The online assessments mix three formats: one-line concept MCQs, numericals meant to
be done by hand in a minute (fit a line to five points, $r^2$ from $r$, the 63\% bootstrap coverage, an IDF, a
split's information gain), and ``which statements are true'' multi-selects (convexity of $\sin x$ and $\abs x$, what
$\lambda$ does, bagging versus boosting, pitfalls of tree importances, $k$-means stopping conditions, why standardise
before PCA). Interviews go deeper and decide the offer: rounds have covered means, distributions, correlation and
covariance, $t$-, $z$- and p-values, confusion matrices, regression and trees; ML-engineer loops ask for derivations
and NumPy implementations with follow-ups several levels deep. Topics absent from those 49 questions are \emph{not}
skipped here; they are simply given less space.

\section*{How each topic is built}
Every topic gets four layers:
\begin{enumerate}
	\ii \textbf{intuition}, with an everyday analogy where one helps;
	\ii \textbf{the precise definition or derivation};
	\ii \textbf{a worked numerical in the shape an OA would set}, computed by code;
	\ii \textbf{the trap} an interviewer or question-setter plants, in a green remark marked \trap
\end{enumerate}
The typography follows Evan Chen's \emph{An Infinitely Large Napkin}: each section opens with a
\emph{prototypical example} in red; new terms appear in \vocab{blue}; theorems sit in blue boxes, worked examples
in red boxes, remarks behind a green bar, quick questions behind a black bar, and a one-line \textbf{moral} in a
green frame closes a topic. Chapter-end problems carry one to three chili peppers for difficulty
(\chili\ easy, \chili\chili\ medium, \chili\chili\chili\ hard); hints and solutions are in the appendices, and every
problem links to both. Chapter~\ref{ch:rapid} is a rapid-fire oral drill; Chapter~\ref{ch:problems} is a set of
OA-style problems.

\section*{How to read it}
The chapters depend on each other roughly as below. Chapter~\ref{ch:maths} is a toolkit to consult rather than
read straight through. For a two-week sprint before an OA, read the starred chapters (setup, linear and logistic
regression, bias--variance, SVM, trees, ensembles, unsupervised, evaluation, data preparation), do the problems,
then drill the rapid-fire chapter aloud.
\begin{center}
	\begin{tikzpicture}[node distance=0.55cm and 0.5cm, font=\scriptsize,
			b/.style={draw, rounded corners, fill=blue!5, minimum width=1.9cm, minimum height=0.55cm, align=center},
			s/.style={b, fill=orange!12}, ->, >=Stealth]
		\node[s] (c1) {1 Setup};
		\node[b, right=of c1] (c2) {2 Maths};
		\node[s, below=of c1] (c3) {3 Linear reg.};
		\node[s, right=of c3] (c4) {4 Bias--variance};
		\node[s, right=of c4] (c5) {5 Logistic};
		\node[b, below=of c3] (c6) {6 Generative};
		\node[b, right=of c6] (c7) {7 $k$-NN};
		\node[s, right=of c7] (c8) {8 SVM};
		\node[s, below=of c6] (c9) {9 Trees};
		\node[s, right=of c9] (c10) {10 Ensembles};
		\node[s, right=of c10] (c11) {11 Unsupervised};
		\node[b, below=of c9] (c12) {12 Semi-sup.};
		\node[s, right=of c12] (c13) {13 Evaluation};
		\node[b, right=of c13] (c14) {14 Statistics};
		\node[s, below=of c12] (c15) {15 Data prep};
		\node[b, right=of c15] (c16) {16 Optimisation};
		\node[b, right=of c16] (c17) {17--19 Practice};
		\draw (c1) -- (c3); \draw (c2) -- (c4); \draw (c3) -- (c4); \draw (c4) -- (c5);
		\draw (c3) -- (c6); \draw (c4) -- (c7); \draw (c5) -- (c8);
		\draw (c6) -- (c9); \draw (c9) -- (c10); \draw (c7) -- (c11);
		\draw (c10) -- (c13); \draw (c13) -- (c14); \draw (c11) -- (c12);
		\draw (c12) -- (c15); \draw (c15) -- (c16); \draw (c16) -- (c17);
	\end{tikzpicture}
\end{center}
{\small Orange: topics that appeared in the OA papers described above.}

\section*{Sources}
The material is standard and is written here from the author's own understanding, then checked numerically. For
deeper reading: James, Witten, Hastie and Tibshirani's \emph{An Introduction to Statistical Learning}
\cite{islr}; Hastie, Tibshirani and Friedman's \emph{The Elements of Statistical Learning} \cite{esl}; Bishop's
\emph{Pattern Recognition and Machine Learning} \cite{bishop}; Murphy's \emph{Probabilistic Machine Learning}
\cite{murphy}; Shalev-Shwartz and Ben-David's \emph{Understanding Machine Learning} \cite{ssbd}; Chen and Guestrin
on XGBoost \cite{xgboost}; Friedman on gradient boosting \cite{friedman2001}. Other papers are cited where a
specific result is used. No benchmark numbers or model sizes are quoted from memory; everything numerical is
computed.

\section*{Reproducibility}
Every worked numerical, table, plotted curve and quoted program output in this book is produced by one of
$\gv{summary/scripts}$ Python scripts in \fpath{code/ml/}, run by \fpath{build.sh}, and inserted into the \LaTeX{}
source automatically --- nothing is typed in by hand. Every from-scratch implementation is checked against the
reference library (scikit-learn, SciPy, XGBoost, LightGBM, PyTorch) with an explicit \code{np.allclose} assertion;
this build passed all $\gv{summary/checks}$ such checks, and a single mismatch stops the build. Seeds are fixed and
everything runs on a CPU in seconds. Versions used for this build: Python $\gv{a00_versions/python}$, NumPy
$\gv{a00_versions/numpy}$, SciPy $\gv{a00_versions/scipy}$, scikit-learn $\gv{a00_versions/sklearn}$, XGBoost
$\gv{a00_versions/xgboost}$, LightGBM $\gv{a00_versions/lightgbm}$, PyTorch \texttt{\gv{a00_versions/torch}},
Matplotlib $\gv{a00_versions/matplotlib}$ (used only to extract contour lines; every figure is drawn with TikZ and
pgfplots). The book contains $\gv{summary/problems}$ problems with hints and solutions. To rebuild: \texttt{./build.sh ml}.

The page layout, theorem boxes and problem/solution machinery are adapted from the source of \emph{An Infinitely
Large Napkin} (GPL v3; see \fpath{latex/NOTICE.md}); none of its text is reused.

\hfill --- Philospher
